\begin{exercises}
  \recommended \item
    下列 \( \nbyn{2} \) 矩阵构成的向量空间中的哪些子集，在继承的运算下是子空间？
    对每一个是子空间的，给出其参数化描述。
    对每一个不是的，给出一条不满足的条件。
    \begin{exparts}
      \partsitem \( \set{\begin{mat}
                      a  &0  \\
                      0  &b
                    \end{mat}  \suchthat a,b\in\Re}  \)
      \partsitem \( \set{\begin{mat}
                      a  &0  \\
                      0  &b
                    \end{mat}  \suchthat a+b=0} \)
      \partsitem \( \set{\begin{mat}
                      a  &0  \\
                      0  &b
                    \end{mat}  \suchthat a+b=5} \)
      \partsitem \( \set{\begin{mat}
                      a  &c  \\
                      0  &b
                    \end{mat}  \suchthat a+b=0, c\in\Re} \)
    \end{exparts}
    \begin{answer}
      根据 \nearbylemma{th:SubspIffClosed}，要判断 $\matspace_{\nbyn{2}}$
      的每个子集是否为子空间，只需检查它是否非空且封闭。
      \begin{exparts}
        \partsitem 是，容易验证它非空且封闭。
          这就是它的一个参数化描述。
          \begin{equation*}
            \set{a\begin{mat}[r]
                    1  &0  \\
                    0  &0
                  \end{mat}
                 +b\begin{mat}[r]
                    0  &0  \\
                    0  &1
                   \end{mat}
                 \suchthat a,b\in\Re}
          \end{equation*}
          顺便指出，这个参数化描述本身也表明它是一个子空间，
          因为它被表示为由两个矩阵构成的集合张成，
          而任何张成都是子空间。
        \partsitem 是；容易验证它非空且封闭。
         或者，如前一个答案所述，参数化描述的存在性
         表明它是一个子空间。
         要得到参数化描述，
         可将条件 $a+b=0$ 改写为 $a=-b$。
         于是有下式。
          \begin{equation*}
            \set{\begin{mat}
                    -b  &0  \\
                    0   &b
                  \end{mat}
                 \suchthat b\in\Re}
            =\set{b\begin{mat}[r]
                    -1  &0  \\
                    0   &1
                  \end{mat}
                 \suchthat b\in\Re}
          \end{equation*}
        \partsitem 不是。
          它对加法不封闭。
          例如，
          \begin{equation*}
            \begin{mat}[r]
              5  &0  \\
              0  &0
            \end{mat}
            +\begin{mat}[r]
              5  &0  \\
              0  &0
            \end{mat}
            =\begin{mat}[r]
              10  &0  \\
              0  &0
            \end{mat}
          \end{equation*}
          不在该集合中。
          （这个集合对标量乘法也不封闭，
          例如，它不包含零矩阵。）
        \partsitem 是。
          \begin{equation*}
            \set{b\begin{mat}[r]
                    -1  &0  \\
                    0   &1
                  \end{mat}
                 +c\begin{mat}[r]
                    0  &1  \\
                    0  &0
                   \end{mat}
                 \suchthat b,c\in\Re}
          \end{equation*}
      \end{exparts}
     \end{answer}
  \recommended \item
    这是 \( \polyspace_2 \) 的子空间吗：
    \( \set{a_0+a_1x+a_2x^2\suchthat a_0+2a_1+a_2=4} \)？
    如果是，给出其参数化描述。
    \begin{answer}
      不是，它不封闭。
      具体来说，它对标量乘法不封闭，因为它
      不包含零多项式。
    \end{answer}
  \item 下面的向量是否在该集合的张成之中？
    \begin{equation*}
      \colvec{1 \\ 0 \\ 3}
      \quad
      \set{\colvec{2 \\ 1 \\ -1},
           \colvec{1 \\ -1 \\ 1}}
    \end{equation*}
    \begin{answer}
      方程
      \begin{equation*}
        \colvec{1 \\ 0 \\ 3} =
                 c_1\colvec{2 \\ 1 \\ -1}
                 +c_2\colvec{1 \\ -1 \\ 1}
      \end{equation*}
      给出一个线性方程组
      \begin{equation*}
        \begin{amat}{2}
          2  &1  &1  \\
          1  &-1 &0  \\
          -1 &1  &3
        \end{amat}
        \grstep[(1/2)\rho_1+\rho_3]{(-1/2)\rho_1+\rho_2}
        \begin{amat}{2}
          2  &1    &1  \\
          0  &-3/2 &-1/2  \\
          0  &0  &3
        \end{amat}
      \end{equation*}
      该方程组无解，所以该向量不在张成之中。
    \end{answer}
  \recommended \item
    判断该向量是否张成于该集合的张成之中，所在
    空间如下。
    \begin{exparts}
      \partsitem \( \colvec[r]{2 \\ 0 \\ 1} \),
        \( \set{\colvec[r]{1 \\ 0 \\ 0},
                \colvec[r]{0 \\ 0 \\ 1}  } \),
        在 \( \Re^3 \) 中
      \partsitem \( x-x^3 \),
        \( \set{x^2,2x+x^2,x+x^3} \),
        在 \( \polyspace_3 \) 中
      \partsitem \( \begin{mat}[r]
                 0  &1  \\
                 4  &2
               \end{mat}  \),
        \( \set{\begin{mat}[r]
                  1  &0  \\
                  1  &1
                \end{mat},
                \begin{mat}[r]
                  2  &0  \\
                  2  &3
                \end{mat}  } \),
        在 \( \matspace_{\nbyn{2}} \) 中
    \end{exparts}
    \begin{answer}
      \begin{exparts}
         \partsitem 是，由
          \begin{equation*}
            r_1\colvec[r]{1 \\ 0 \\ 0}+r_2\colvec[r]{0 \\ 0 \\ 1}
              =\colvec[r]{2 \\ 0 \\ 1}
          \end{equation*}
          得到的线性方程组的解给出 \( r_1=2 \) 与 \( r_2=1 \)。
        \partsitem 是；由
          \( r_1(x^2)+r_2(2x+x^2)+r_3(x+x^3)=x-x^3 \)
          \begin{equation*}
            \begin{linsys}{3}
                  &  &2r_2 &+ &r_3 &= &1  \\
              r_1 &+ &r_2  &  &    &= &0  \\
                  &  &     &  &r_3 &= &-1
            \end{linsys}
          \end{equation*}
          可得 \( -1(x^2)+1(2x+x^2)-1(x+x^3)=x-x^3 \)。
       \partsitem 不是；所给两个矩阵的任意组合，其右上角元素都为零。
      \end{exparts}
    \end{answer}
  \item % \cite{Cleary} credit removed at request Cleary 2022-Sep-13
    一位超级英雄位于二维平面上的原点处。
    超级英雄有两件装置：一块悬浮滑板，可沿方向 $\binom{3}{1}$
    移动任意距离；
    以及一张魔毯，可沿方向
    $\binom{1}{2}$ 移动任意距离。
    \begin{exparts}
      \partsitem 一个邪恶反派藏身于平面上点
        $(-5,7)$ 处。
        超级英雄要移动多少个悬浮滑板单位和多少个魔毯单位，
        才能到达反派所在处？
     \partsitem 平面上是否存在某个地方，可以让反派安全藏身而不被找到？
        如果存在，给出一个这样的位置。
        如果不存在，解释原因。
     \partsitem 超级英雄和反派被传送到了一个三维空间，此时超级英雄拥有三件装置。
        \begin{equation*}
          \text{悬浮滑板: }\colvec{-1 \\ 0 \\ 3}
          \quad
          \text{魔毯: }\colvec{2 \\ 1 \\ 0}
          \quad
          \text{滑板车: }\colvec{5 \\ 4 \\ -9}
        \end{equation*}
        是否还有可以让反派安全藏身的地方？
        如果有，给出一个这样的位置；如果没有，解释原因。
    \end{exparts}
    \begin{answer}
      \begin{exparts}
        \partsitem 关系式
          \begin{equation*}
            \colvec{-5 \\ 7}=a\colvec{3 \\ 1}+b\colvec{1 \\ 2}
          \end{equation*}
          给出如下方程组。
          \begin{equation*}
            \begin{linsys}{2}
              3a &+ &b  &= &-5 \\
               a &+ &2b &= &7
            \end{linsys}
            \grstep{-(1/3)\rho_1+\rho_2}
            \begin{linsys}{2}
              3a &+ &b      &= &-5 \\
                 &  &(5/3)b &= &26/3
            \end{linsys}
          \end{equation*}
          所以 $b=26/5$。
          代入 $3a+(26/5)=-5$ 得 $a=-17/5$。
          用本小节的术语来说，向量 $\binom{-5}{7}$
          在另外两个向量构成的集合的张成之中。
        \partsitem 没有藏身之处。
          下面的高斯消元法化简
          \begin{equation*}
            \begin{linsys}{2}
              3a &+ &b  &= &x \\
               a &+ &2b &= &y
            \end{linsys}
            \grstep{-(1/3)\rho_1+\rho_2}
            \begin{linsys}{2}
              3a &+ &b      &= &x\hfill\hbox{} \\
                 &  &(5/3)b &= &-(1/3)x+y
            \end{linsys}
          \end{equation*}
          表明，对任意 $x,y\in\R$ 都存在一对 $a,b\in\R$。
          用本小节的术语来说，这两个向量的张成
          是整个平面。
        \partsitem 没有藏身之处。
          考虑这个方程组。
          \begin{equation*}
            a\colvec{-1 \\ 0 \\ 3}+b\colvec{2 \\ 1 \\ 0}+c\colvec{5 \\ 4 \\ -9}
                = \colvec{x \\ y \\ z}
          \end{equation*}
          化简后
          \begin{align*}
            \begin{linsys}{3}
              -a &+ &2b &+ &5c &= &x \\
                 &  &b  &+ &4c &= &y \\
              3a &  &   &- &9c &= &z \\
            \end{linsys}
            & \grstep{3\rho_1+\rho_3}
            \begin{linsys}{3}
              -a &+ &2b  &+ &5c &= &x\hfill\hbox{} \\
                 &  &b   &+ &4c &= &y\hfill\hbox{} \\
                 &  &6b  &+ &6c &= &3x+z \\
            \end{linsys}                                  \\
            & \grstep{-6\rho_2+\rho_3}
            \begin{linsys}{3}
              -a &+ &2b  &+ &5c  &= &x\hfill\hbox{} \\
                 &  &b   &+ &4c  &= &y\hfill\hbox{} \\
                 &  &    & &-18c &= &3x+-6y+z \\
            \end{linsys}
          \end{align*}
          表明，对任意 $x,y,z\in\R$，
          所需的 $a$、$b$ 与~$c$ 都存在。
          超级英雄的装置张成了整个空间。
      \end{exparts}
    \end{answer}
  \item
    下列哪些是实变量实值函数构成的向量空间中
    张成
    \( \spanof{\set{\cos^2x,\sin^2x} } \)
    的成员？
    \begin{exparts*}
      \partsitem \( f(x)=1 \)
      \partsitem \( f(x)=3+x^2 \)
      \partsitem \( f(x)=\sin x \)
      \partsitem \( f(x)=\cos (2x) \)
    \end{exparts*}
    \begin{answer}
      \begin{exparts}
        \partsitem 是；它在那个张成之中，因为
          \( 1\cdot\cos^2x+1\cdot\sin^2x=f(x) \)。
        \partsitem 不是，因为 \( r_1\cos^2x+r_2\sin^2x=3+x^2 \) 没有对
          所有 \( x \) 都成立的标量解。
          例如，取 $x$ 为 $0$ 和 $\pi$ 得到两个方程 $r_1\cdot 1+r_2\cdot 0=3$ 与
          $r_1\cdot 1+r_2\cdot 0=3+\pi^2$，二者
          相互矛盾。
        \partsitem 不是；考虑取 $x$ 为 $\pi/2$ 与
          $3\pi/2$ 时的结果。
        \partsitem 是，\( \cos (2x)=1\cdot\cos^2(x)-1\cdot\sin^2(x) \)。
      \end{exparts}
     \end{answer}
  \recommended \item
    下列哪些集合张成 \( \Re^3 \)？
    也就是说，下列哪些集合具有这样的性质：任意三元
    向量都可以表示为该集合元素的某个适当
    线性组合？
    \begin{exparts*}
      \partsitem \( \set{ \colvec[r]{1 \\ 0 \\ 0},
               \colvec[r]{0 \\ 2 \\ 0},
               \colvec[r]{0 \\ 0 \\ 3}  } \)
      \partsitem \( \set{ \colvec[r]{2 \\ 0 \\ 1},
               \colvec[r]{1 \\ 1 \\ 0},
               \colvec[r]{0 \\ 0 \\ 1}  } \)
      \partsitem \( \set{ \colvec[r]{1 \\ 1 \\ 0},
               \colvec[r]{3 \\ 0 \\ 0}  } \)
      \partsitem \( \set{ \colvec[r]{1 \\ 0 \\ 1},
               \colvec[r]{3 \\ 1 \\ 0},
               \colvec[r]{-1 \\ 0 \\ 0},
               \colvec[r]{2 \\ 1 \\ 5}  } \)
      \partsitem \( \set{ \colvec[r]{2 \\ 1 \\ 1},
               \colvec[r]{3 \\ 0 \\ 1},
               \colvec[r]{5 \\ 1 \\ 2},
               \colvec[r]{6 \\ 0 \\ 2}  } \)
    \end{exparts*}
    \begin{answer}
      \begin{exparts}
         \partsitem 是，对任意 \( x,y,z\in\Re \)，方程
          \begin{equation*}
             r_1\colvec[r]{1 \\ 0 \\ 0}
             +r_2\colvec[r]{0 \\ 2 \\ 0}
             +r_3\colvec[r]{0 \\ 0 \\ 3}
             =\colvec{x \\ y \\ z}
          \end{equation*}
          有解 \( r_1=x \)，\( r_2=y/2 \) 与
          \( r_3=z/3 \)。
        \partsitem 是，方程
          \begin{equation*}
            r_1\colvec[r]{2 \\ 0 \\ 1}
            +r_2\colvec[r]{1 \\ 1 \\ 0}
            +r_3\colvec[r]{0 \\ 0 \\ 1}
            =\colvec{x \\ y \\ z}
          \end{equation*}
          给出如下
          \begin{equation*}
            \begin{linsys}{3}
              2r_1 &+  &r_2  &  &    &=  &x \\
                   &   &r_2  &  &    &=  &y \\
               r_1 &   &     &+ &r_3 &=  &z \\
            \end{linsys}
            \grstep{-(1/2)\rho_1+\rho_3}\repeatedgrstep{(1/2)\rho_2+\rho_3}
            \begin{linsys}{3}
              2r_1 &+  &r_2  &  &    &=  &x\hfill\hbox{} \\
                   &   &r_2  &  &    &=  &y\hfill\hbox{} \\
                   &   &     &  &r_3 &=  &-(1/2)x+(1/2)y+z \\
            \end{linsys}
          \end{equation*}
          于是，给定任意 $x$、$y$ 和 $z$，我们可以算出
          \( r_3=(-1/2)x+(1/2)y+z \)、\( r_2=y \) 与
          \( r_1=(1/2)x-(1/2)y \)。
        \partsitem 不是。
          特别地，我们无法得到向量
          \begin{equation*}
            \colvec[r]{0 \\ 0 \\ 1}
          \end{equation*}
          的线性组合，因为所给的两个
          向量的第三个分量都为零。
       \partsitem 是。
         方程
         \begin{equation*}
           r_1\colvec[r]{1 \\ 0 \\ 1}
           +r_2\colvec[r]{3 \\ 1 \\ 0}
           +r_3\colvec[r]{-1\\ 0 \\ 0}
           +r_4\colvec[r]{2 \\ 1 \\ 5}
           =\colvec{x \\ y \\ z}
         \end{equation*}
         导出如下化简。
         \begin{equation*}
           \begin{amat}{4}
             1  &3  &-1  &2  &x  \\
             0  &1  &0   &1  &y  \\
             1  &0  &0   &5  &z
           \end{amat}
           \grstep{-\rho_1+\rho_3}\repeatedgrstep{3\rho_2+\rho_3}
           \begin{amat}{4}
             1  &3  &-1  &2  &x \\
             0  &1  &0   &1  &y  \\
             0  &0  &1   &6  &-x+3y+z
           \end{amat}
         \end{equation*}
         它有无穷多解。
         例如，我们可以取 $r_4$ 为零，然后按通常的
         回代方法，用 $x$、$y$ 和 $z$ 解出
         $r_3$、$r_2$ 和 $r_1$。
       \partsitem 不是。
         方程
         \begin{equation*}
           r_1\colvec[r]{2 \\ 1 \\ 1}
           +r_2\colvec[r]{3 \\ 0 \\ 1}
           +r_3\colvec[r]{5 \\ 1 \\ 2}
           +r_4\colvec[r]{6 \\ 0 \\ 2}
           =\colvec{x \\ y \\ z}
         \end{equation*}
         导出如下化简。
         \begin{multline*}
           \begin{amat}{4}
             2  &3  &5   &6  &x  \\
             1  &0  &1   &0  &y  \\
             1  &1  &2   &2  &z
           \end{amat}                                             \\
           \grstep[-(1/2)\rho_1+\rho_3]{-(1/2)\rho_1+\rho_2}
           \repeatedgrstep{-(1/3)\rho_2+\rho_3}
           \begin{amat}{4}
             2  &3     &5     &6  &x \\
             0  &-3/2  &-3/2  &-3 &-(1/2)x+y  \\
             0  &0     &0     &0  &-(1/3)x-(1/3)y+z
           \end{amat}
         \end{multline*}
         这表明并非每个三元向量都能如此表示。
         只有满足约束 $-(1/3)x-(1/3)y+z=0$ 的向量才在张成之中。
         （要看出任何这样的向量确实可以表示，
         可取 $r_3$ 与 $r_4$
         为零，然后用回代法以 $x$、$y$ 和
         $z$ 解出 $r_1$ 与 $r_2$。）
      \end{exparts}
    \end{answer}
  \recommended \item
    对每个子空间的描述给出参数化。
    然后把每个子空间表示为一个张成。
    \begin{exparts}
      \partsitem  三宽行向量中的子集
        \( \set{\rowvec{a &b &c}\suchthat a-c=0}   \)
      \partsitem \( \matspace_{\nbyn{2}} \) 的这个子集
        \begin{equation*}
          \set{\begin{mat}
                 a  &b  \\
                 c  &d
               \end{mat}  \suchthat a+d=0}
        \end{equation*}
      \partsitem \( \matspace_{\nbyn{2}} \) 的这个子集
        \begin{equation*}
          \set{\begin{mat}
                 a  &b  \\
                 c  &d
               \end{mat}  \suchthat \text{\( 2a-c-d=0 \)
                                               且 \( a+3b=0 \)} }
        \end{equation*}
      \partsitem 子集 \( \set{a+bx+cx^3\suchthat a-2b+c=0} \)，
        属于 \( \polyspace_3 \)
      \partsitem \( \polyspace_2 \) 的子集：满足 \( p(7)=0 \) 的
        次数不超过二的多项式 \( p \)
    \end{exparts}
    \begin{answer}
      \begin{exparts}
        \partsitem \( \set{\rowvec{c &b &c}\suchthat b,c\in\Re}
                      =\set{b\rowvec{0 &1 &0}+c\rowvec{1 &0 &1}
                        \suchthat b,c\in\Re} \)
           张成集合的显然选择是
           $\set{\rowvec{0 &1 &0},\rowvec{1 &0 &1}}$。
        \partsitem \( \set{\begin{mat}
                       -d &b  \\
                       c  &d
                      \end{mat} \suchthat b,c,d\in\Re}
                     =\set{b\begin{mat}[r]
                       0  &1  \\
                       0  &0
                      \end{mat}
                     +c\begin{mat}[r]
                       0  &0  \\
                       1  &0
                      \end{mat}
                     +d\begin{mat}[r]
                       -1  &0  \\
                       0  &1
                      \end{mat}  \suchthat b,c,d\in\Re} \)
            张成这个空间的一个集合就由这三个矩阵组成。
        \partsitem 方程组
          \begin{equation*}
            \begin{linsys}{4}
              a  &+  &3b  &   &   &  &  &=  &0  \\
             2a  &   &    &   &-c &- &d &=  &0
            \end{linsys}
          \end{equation*}
          给出 \( b=-(c+d)/6 \) 与 \( a=(c+d)/2 \)。
          于是一种描述如下。
          \begin{equation*}
            \set{c\begin{mat}[r]
                       1/2  &-1/6  \\
                       1    &0
                      \end{mat}
                     +d\begin{mat}[r]
                       1/2  &-1/6  \\
                       0    &1
                      \end{mat}  \suchthat c,d\in\Re}
           \end{equation*}
          这表明张成该子空间的一个集合由那
          两个矩阵组成。
        \partsitem 由 $a=2b-c$ 可知，集合
           \( \set{(2b-c)+bx+cx^3 \suchthat b,c\in\Re} \)
           等于集合
           \( \set{b(2+x)+c(-1+x^3) \suchthat b,c\in\Re}  \)。
           所以该子空间是集合 $\set{2+x, -1+x^3}$ 的张成。
        \partsitem 集合
          \( \set{a+bx+cx^2\suchthat a+7b+49c=0} \)
          可以参数化为
          \begin{equation*}
            \set{b(-7+x)+c(-49+x^2)\suchthat b,c\in\Re}
          \end{equation*}
          因此
          有张成集 $\set{-7+x,-49+x^2}$。
      \end{exparts}
    \end{answer}
  \recommended \item
    对给定空间中的给定子空间，求一个张成它的集合。
    （\textit{提示。}先对每一个作参数化。）
    \begin{exparts}
      \partsitem  \( \Re^3 \) 中的 \( xz \) 平面
      \partsitem \( \set{\colvec{x \\ y \\ z}\suchthat 3x+2y+z=0} \)
            在 \( \Re^3 \) 中
      \partsitem \( \set{\colvec{x \\ y \\ z \\ w}\suchthat
                       2x+y+w=0 \text{\ 且\ } y+2z=0} \)
            在 \( \Re^4 \) 中
      \partsitem \( \set{a_0+a_1x+a_2x^2+a_3x^3\suchthat
                        a_0+a_1=0 \text{\ 且\ } a_2-a_3=0} \)
            在 \( \polyspace_3 \) 中
      \partsitem 空间 \( \polyspace_4 \) 中的集合 \( \polyspace_4 \)
      \partsitem \( \matspace_{\nbyn{2}} \) 中的 \( \matspace_{\nbyn{2}} \)
    \end{exparts}
    \begin{answer}
      \begin{exparts}
        \partsitem 我们可以这样参数化
          \begin{equation*}
             \set{\colvec{x \\ 0 \\ z}\suchthat x,z\in\Re}
             =\set{x\colvec[r]{1 \\ 0 \\ 0}
                  +z\colvec[r]{0 \\ 0 \\ 1}\suchthat x,z\in\Re}
          \end{equation*}
          由此得到张成集如下。
          \begin{equation*}
             \set{\colvec[r]{1 \\ 0 \\ 0},\colvec[r]{0 \\ 0 \\ 1}}
          \end{equation*}
        \item 这是参数化描述，以及相应的张成集。
          \begin{equation*}
             \set{y\colvec[r]{-2/3 \\ 1 \\ 0}+z\colvec[r]{-1/3 \\ 0 \\ 1}
                 \suchthat y,z\in\Re }
             \qquad
             \set{\colvec[r]{-2/3 \\ 1 \\ 0},\colvec[r]{-1/3 \\ 0 \\ 1} }
          \end{equation*}
        \partsitem \( \set{\colvec[r]{1 \\ -2 \\ 1 \\ 0},
                      \colvec[r]{-1/2 \\ 0 \\ 0 \\ 1} } \)
        \partsitem 把描述参数化为
          \( \set{-a_1+a_1x+a_3x^2+a_3x^3\suchthat a_1,a_3\in\Re } \)
          得到 \( \set{-1+x,x^2+x^3}. \)
        \partsitem \( \set{1,x,x^2,x^3,x^4} \)
        \partsitem \( \set{ \begin{mat}[r]
                   1  &0  \\
                   0  &0
                 \end{mat},
                 \begin{mat}[r]
                   0  &1  \\
                   0  &0
                 \end{mat},
                 \begin{mat}[r]
                   0  &0  \\
                   1  &0
                 \end{mat},
                 \begin{mat}[r]
                   0  &0  \\
                   0  &1
                 \end{mat} } \)
      \end{exparts}
    \end{answer}
  \item
    \( \Re^2 \) 是 \( \Re^3 \) 的子空间吗？
    \begin{answer}
      严格地说，不是。
      \( \Re^3 \) 的子空间都是三元向量构成的集合，而
      \( \Re^2 \) 是二元向量构成的集合。
      不过显然，\( \Re^2 \) 与 \( \Re^3 \) 的这个子空间“几乎
      一样”。
      \begin{equation*}
        \set{\colvec{x \\ y \\ 0}\suchthat x,y\in\Re}
      \end{equation*}
    \end{answer}
  \recommended \item
    判断下列每一个是否为实变量实值函数构成的向量空间的子空间。
    \begin{exparts}
      \partsitem \definend{偶}\index{function!even}\index{even functions} 函数
        \( \set{\map{f}{\Re}{\Re} \suchthat f(-x)=f(x) \text{对一切 } x} \)。
        例如，该集合的两个成员是 $f_1(x)=x^2$
        与 $f_2(x)=\cos (x)$。
      \partsitem \definend{奇}\index{function!odd}\index{odd function}
        函数
        \( \set{\map{f}{\Re}{\Re} \suchthat f(-x)=-f(x) \text{对一切 } x} \)。
        两个成员是 $f_3(x)=x^3$ 与 $f_4(x)=\sin(x)$。
    \end{exparts}
    \begin{answer}
      当然，加法与标量乘法运算就是
      从外围空间继承的那两个运算。
      \begin{exparts}
        \partsitem 这是一个子空间。
          它不是空集，因为它至少包含上面给出的两个示例函数。
          它是封闭的，因为若 \( f_1,f_2 \) 是偶函数且
          \( c_1,c_2 \) 是标量，则有下式。
          \begin{equation*}
            (c_1f_1+c_2f_2)\,(-x)
            =c_1\,f_1(-x)+c_2\,f_2(-x)
            =c_1\,f_1(x)+c_2\,f_2(x)
            =(c_1f_1+c_2f_2)\,(x)
          \end{equation*}
        \partsitem 这也是一个子空间；验证方法与
          前一个类似。
      \end{exparts}
    \end{answer}
  \item
    \nearbyexample{ex:SpanSingVec} 说明，对向量空间 $V$ 中任一向量 $\vec{v}$，
    集合 $\set{r\cdot\vec{v}\suchthat r\in\Re}$
    是 $V$ 的一个子空间。
    （这不过是单元素集 $\set{\vec{v}}$ 的
    张成\index{span!of a singleton}。）
    这样的子空间是否一定是真子空间？
    \begin{answer}
      它可以不是真子空间。
      例如，
      若向量空间是 \( \Re^1 \) 且 \( \vec{v}\neq\zero\, \)，
      则子空间 $\set{r\cdot\vec{v}\suchthat r\in\Re}$
      就是整个 $\Re^1$。
    \end{answer}
  \item
    向量空间定义之后的一个例子表明，
    齐次线性方程组的解集是一个向量空间。
    用本小节的术语来说，当方程组有 $n$ 个变量时，它是 $\Re^n$ 的一个子空间。
    那么非齐次线性方程组呢；它的解是否构成一个
    子空间（在继承的运算下）？
    \begin{answer}
      不构成，这样的集合不封闭。
      首先，它不包含零向量。
    \end{answer}
  \item \cite{Cleary}
   就下列每一条给出一个例子，或者解释为什么
   不可能给出。
   \begin{exparts}
     \item $\matspace_{\nbyn{2}}$ 的一个不是
       子空间的非空子集。
     \item $\Re^2$ 中一个不张成该空间的两向量集合。
   \end{exparts}
   \begin{answer}
     \begin{exparts}
       \item $\matspace_{\nbyn{2}}$ 的这个非空子集不是子空间。
         \begin{equation*}
           A=\set{
             \begin{mat}[r]
               1 &2 \\
               3 &4
             \end{mat},
             \begin{mat}[r]
               5 &6 \\
               7 &8
             \end{mat}}
         \end{equation*}
         它不是 $\matspace_{\nbyn{2}}$ 的子空间的一个理由是
         它不包含零矩阵。
         （另一个理由是它对加法不封闭，因为两者的和
         不是 $A$ 的元素。
         它对标量乘法也不封闭。）
        \item 这个由两个向量构成的集合不张成 $\Re^2$。
          \begin{equation*}
            \set{\colvec[r]{1 \\ 1},\colvec[r]{3 \\ 3}}
          \end{equation*}
          这两个向量的任何线性组合都无法给出
          第二个分量不等于第一个分量的向量。
      \end{exparts}
   \end{answer}
  \item
    \nearbyexample{ex:SubspRThree} 表明
    $\Re^3$ 有无穷多个子空间。
    是否每个非平凡空间都有无穷多个子空间？
    \begin{answer}
      不是。
      \( \Re^1 \) 仅有的子空间是它本身与其
      平凡子空间。
      $\Re$ 的任何包含非零成员 $\vec{v}$ 的子空间 $S$
      必须包含其全部标量倍数构成的集合
      $\set{r\cdot\vec{v}\suchthat r\in\Re}$。
      但这个集合就是整个 $\Re$。
    \end{answer}
  \item \label{exer:SubspIffClosed}
    补全 \nearbylemma{th:SubspIffClosed} 的证明。
    \begin{answer}
      第~(1) 项已在正文中验证。

